\section{Summary}
\textbf{RQ1 -- privacy.} Phương pháp nối REM, noisy reuse và cap thành cận không gian cho lịch sử đã bảo vệ;
kernel chọn đoạn có common support và cận riêng cho toàn transcript. Hậu xử lý dẫn tới cận minimum giữa
hai nhánh, cùng diễn giải posterior odds. Kết quả lý thuyết phân biệt kernel lý tưởng với mã hiện tại:
float sampler Geo-I và full upstream restore chưa được chứng nhận; fresh uniform bits và trusted ledger
vẫn là giả thiết của cận Q. Năm Q hoặc Z không công bố không tự tạo anonymity cho người/xe.

\textbf{RQ2 -- utility.} Private adaptation tăng kỳ vọng score lập kế hoạch; tăng actual utility cần proxy phù hợp.
Union/cache và ranking bảo toàn POI đã có dưới cùng GPS/request/order. Cận theo matched public-base
và cận chuyển có điều kiện không được phát biểu thành floor phổ quát. Thực nghiệm hiện tại đạt
\PrimaryRecall\% Recall với CI \PrimaryCI\%, giảm khoảng 0,4 pp so với ba adaptation gần 100\%.

\textbf{RQ3 -- đánh đổi và đối chiếu.} Trên cohort đã xem, GeoI-BQ có Hit100 thấp hơn ba adaptation trong
bank kiểm thử cho S1--S3/S9--S10. Chi phí 155 point queries/phiên cao hơn singleton controls, nhưng không
cần fake packets phụ. Benchmark được tổ chức theo output contract, bảo lưu metrics gốc khi dùng được,
N/A, failures và CIs; các kết quả không là xác nhận vượt mọi implementation chính thức của paper.

Đóng góp thực sự là đặc tả bài toán/đánh giá, thiết kế kênh Q có cận toàn phiên và phép đối chiếu tái lập,
không phải một metric MAE/Recall mới hay tuyên bố bảo vệ hoàn toàn mọi scenario.
Các bước tiếp theo gồm chứng nhận sampler/restore, phân tích nhiều phiên và side channels,
đánh giá độc lập trên thành phố/catalogue động, tái lập predictor gốc và calibration utility thực địa.
S4--S6/S8 cần mô hình target/dữ liệu bổ sung khi trở thành thực nghiệm chính. Các hướng này không
thay đổi kết luận có điều kiện và bằng chứng hiện đã được báo cáo.
